\section{Measures and inference}\label{sec:method}

\subsection{Measures}

\paragraph{The settlement price.} The settlement VWAP is computed from the window's trades.
It is then compared against the mid at the window's open (the \emph{drift}), against the
closing price (the \emph{reversal}), and against the final minute's own VWAP (the
\emph{terminal gap}). The reversal is the most informative of the three: a price that is
pushed and then released leaves the close somewhere other than the settlement.

Tests use absolute values. A participant who is long the future wants the settlement high and
one who is short wants it low, so across many securities and months the signed drift should
average to roughly nothing whether or not anyone is pushing. The dispersion carries the
information.

\paragraph{Pushing or churning.} The variance ratio compares the variance of a
\VrHorizonSeconds{}-second mid return against \VrHorizonSeconds{} times the variance of a
one-second return. Under a random walk it is one. Sustained pressure in one direction pushes
it above one; heavier two-sided trading, which arrives as bid-ask bounce, pushes it below.
The first-order autocorrelation of signed volume per second asks the same question of order
flow. Impact is split into the part that survives: for each second carrying substantial
one-sided flow, the mid move after ten seconds and after sixty is measured, and the ratio is
the permanent share.

Impact is measured on the mid rather than on trade prices. Signing trades by the tick rule
and then measuring their own subsequent return partly measures the bid-ask bounce reverting,
which makes the estimated impact of a buy come out negative and the permanent share a ratio
of two negative numbers.

\paragraph{Concealment.} An order is treated as concealing size when its displayed quantity
is positive and below its total; a displayed quantity of zero means no instruction was given,
which is the ordinary case. From the book we take the share of resting size that is
concealed, and the replenishment count per thousand shares matched, which is the rate at
which concealed size is converted into executions.

\paragraph{The cost of moving the price.} For each snapshot the displayed book is swept from
the touch by a fixed notional and the resulting price move recorded. Sweeping a fixed size and
measuring the move, rather than asking what a fixed move costs, avoids censoring: the book
carries a finite number of levels, and on a liquid security a ten-basis-point move lies beyond
all of them. The figure is a lower bound on what moving the settlement would take, since the
sweep assumes the book neither replenishes nor reacts, and should be read as an order of
magnitude and as a comparison across groups.

\paragraph{Order behaviour.} The cancel-to-entry ratio, the share of orders cancelled within
one second of entry, the immediate-or-cancel rate (reported separately for the final five
minutes), spreads, displayed depth, and the concentration of volume across the window's
minutes.

\subsection{Paired tests}

Every hypothesis has the same form: a per-security quantity on an expiry session against the
same quantity on that month's control session. Each results file is first reduced to one value
per security and session, so a security contributes one matched pair per month.

Tests are paired $t$-tests with a Wilcoxon signed-rank test alongside, since these
distributions are heavy-tailed. Effect sizes are Cohen's $d$ on the paired differences.
\TestSpecified{} hypotheses are tested over \TestPairsTotal{} matched pairs, and significance
is assessed under Benjamini--Hochberg control of the false discovery rate at
$\alpha = \TestAlpha{}$.

Rejection and support are reported separately. The paired test is two-sided, so a small
$p$-value says a quantity differs between expiry and control sessions---not that it differs in
the direction the hypothesis claims. Each hypothesis therefore carries the direction it
states, and a rejection running the other way is reported as a refutation rather than a
confirmation. Four of them are.

\subsection{Difference-in-differences}

The paired tests establish that a quantity differs on expiry Thursdays. They do not establish
that the settlement mechanism is the reason, and this is where the three-group design in
Section~\ref{sec:design} does its work.

Both contrasts are computed on per-security expiry-minus-control differences, so each
security contributes one observation and the comparison is over securities rather than over
security-months. The twelve months of one security are not twelve independent observations of
how that security responds, and treating them as such would overstate the precision by a
factor of more than three.

The first contrast compares the derivatives groups against the placebo group. If the placebo
securities move the same way on expiry Thursdays, the effect belongs to the calendar rather
than to the settlement, and the paired results above describe the last Thursday of the month
rather than the settlement mechanism.

The second compares the illiquid group against the liquid group. If the settlement incentive
is being acted on, the effect should be larger where moving the price is cheaper---a
prediction the settlement story makes and the calendar explanations mostly do not, since
index rebalancing and month-end flows have no particular reason to scale inversely with
liquidity.

Both are reported in Table~\ref{tab:contrasts}.
